%------------------------------------------------------------------------------
% 4.2 تصميم منظومة الوكلاء
% المالك: AI (امتداد مهمة ٢.٢.١ — قد يستدعي إعادة توزيع داخل الفريق)
% الحجم المستهدف: 3-4 صفحات
% الغرض: التصميم التنفيذي لما نمذجناه تحليلياً في القسم 3.5 — مستوى العقود
%   والأنواع وقواعد التوجيه، لا مستوى الوصف الوظيفي.
%
% تحذيرات:
%   - لا إعادة رسم للشكل 3-7 (fig:agent-state-graph) ولا الشكل 3-8
%     (fig:agent-sequence): أحِل إليهما بـ \ref{} وابنِ فوقهما
%   - لا تكرار لجدول الحقول التحليلي في 3.5.2: هنا المستوى تصميمي — الأنواع
%     (TypedDict)، القيم الابتدائية، قواعد التحقق، ودوال الاختزال (Reducers)
%
% عناصر مطلوبة:
%   [ ] 4.2.1 طوبولوجيا الرسم البياني ومواصفات العقود (tab:agent-contracts):
%       لكل عقدة — التوقيع (مدخل/مخرج)، السلوك عند الفشل، المهلة القصوى،
%       وسياسة إعادة المحاولة
%   [ ] 4.2.2 تصميم الحالة المشتركة: بنية AgentState التنفيذية وعلاقتها بسجل
%       الجلسة المعمّر في قاعدة بيانات النظام (ما يبقى في الذاكرة وما يُثبَّت)
%   [ ] 4.2.3 منطق الحواف المشروطة: جدول قواعد التوجيه (tab:edge-rules) —
%       الحالة / الشرط / الوجهة التالية؛ يشمل حلقة التصحيح الذاتي لوكيل SQL
%       (حد أقصى 3 محاولات مع تغذية Few-Shot من رسالة الخطأ) — اربطها بمتطلب
%       المتانة المعني من جدول المتطلبات الوظيفية
%   [ ] 4.2.4 تصميم الأدوات (tab:tool-schemas): منفذ SQL للقراءة فقط، مسترجع
%       Qdrant، مصنف النوايا، كاشف القيم الشاذة، مولد التقارير — لكل أداة:
%       المدخل / المخرج / القيود الأمنية
%
% مرجع الدقة: الرسم البياني والأدوات الفعلية في مستودع Bayan
%   (backend/src/agents) — صمم من الكود الحقيقي لا من التصور
%------------------------------------------------------------------------------
\section{تصميم منظومة الوكلاء}
\label{sec:agents-design}
~

\subsection{طوبولوجيا الرسم البياني ومواصفات العقود}
\label{subsec:agents-topology}
~

\subsection{تصميم الحالة المشتركة}
\label{subsec:agents-state}
~

\subsection{منطق الحواف المشروطة وقواعد التوجيه}
\label{subsec:agents-edges}
~

\subsection{تصميم الأدوات}
\label{subsec:agents-tools}
~
